% =============================================================================
% Methods 部分精确替换/插入指南
% =============================================================================

% ═══════════════════════════════════════════════════════════════════════════════
% 操作 1 [R]: 替换 Section 2.1 中 "For the present study..." 整段
% 原文位置: 约 line 187-194
% 查找关键字: "Qwen API", "OpenClaw-based runtime", "large language model"
% ═══════════════════════════════════════════════════════════════════════════════

For the present study, \dentalclaw{} was deployed on a server equipped with an
NVIDIA GeForce RTX 3090 GPU. The agent decision layer used the OpenClaw~\cite{openclaw}
runtime with DeepSeek~V4~Pro as the primary reasoning backend and a
direct DeepSeek~API fallback for robustness. Academic literature search drew from
two free, API-key-free sources: Europe~PMC (biomedical literature,
\textgreater1,200~hits for dental~AI queries) and the arXiv~API (computer science
preprints). A DuckDuckGo fallback was retained for non-academic resources.

The agent was responsible for three tasks: (1)~interpreting user research
instructions in the context of the platform's registered methods, codebase assets,
and academic search results; (2)~selecting the most appropriate workflow from the
method registry, or proposing an external solution with a confidence score when no
registered method matched; and (3)~generating an executable configuration plan
(\texttt{platform\_plan.json}) that recorded the decision rationale, selected
method, cited references, and execution prerequisites. The deterministic
workflow execution, data processing, training, evaluation, and reporting steps
remained unchanged: they were carried out by versioned, replayable modules with
audit logs retained for reproducibility.


% ═══════════════════════════════════════════════════════════════════════════════
% 操作 2 [+]: 在 Section 2.1 末尾新增子节
% 插入位置: 原文 2.1 最后一段之后，2.2 标题之前
% ═══════════════════════════════════════════════════════════════════════════════

\paragraph{Agent decision pipeline.}
The agent decision pipeline operated in five stages:

\begin{enumerate}
    \item \textbf{Intent Parsing.} The user's natural-language input was parsed
    into structured fields: \texttt{dataset}, \texttt{task\_family},
    \texttt{modality}, \texttt{mode}, and optional TTA/ensemble flags. A
    bilingual detector (Chinese/English) selected the interface language
    automatically. If both \texttt{dataset} and \texttt{task\_family} were
    unrecognised, the platform returned a polite out-of-scope rejection message
    in the detected language.

    \item \textbf{Academic Literature Search.} A task-aware query was constructed
    from the parsed intent and sent to Europe~PMC and arXiv. Results---paper
    titles, DOIs, and publication years---were aggregated into a structured list
    for the agent.

    \item \textbf{Registry Lookup.} The offline method registry (six pre-validated
    executable routes) was queried for a matching method by task family, modality,
    dataset, and mode. Partial matches and status flags were recorded.

    \item \textbf{Agent Decision.} The agent received the parsed intent, search
    results, registry match status, and a description of available codebase assets
    (nnU-Net scripts, YOLO detection modules, auto-training skills, etc.).
    It produced a JSON decision with one of three outcomes:
    \texttt{use\_registry} (use the matched registry method),
    \texttt{external\_proposal} (propose a solution from web search or codebase
    knowledge, with a confidence score and entrypoint),
    or \texttt{reject} (no viable route).

    \item \textbf{Plan Assembly.} The final \texttt{platform\_plan.json} recorded
    the selected method, decision rationale, cited references, execution
    prerequisites, and whether human review was required before execution
    (\texttt{external\_suggestion} status).
\end{enumerate}

The key design principle was that \textit{external proposals are never executed
automatically}. Methods proposed by the agent with \texttt{external\_suggestion}
status require an explicit \texttt{---allow-external} flag from the user. This
ensured that the platform's six pre-validated registry routes remained the only
automatically executable paths, while the agent could still surface promising
codebase assets for human review.


% ═══════════════════════════════════════════════════════════════════════════════
% 操作 3 [+]: 在 Section 2.2 末尾新增 TTA/Ensemble 小段
% 插入位置: 原文 "2.2 Datasets and representative workflows" 的末尾
% ═══════════════════════════════════════════════════════════════════════════════

In addition to single-model inference, the platform supported test-time
augmentation (TTA) and multi-model ensemble inference through a dedicated module
(\texttt{tta\_ensemble.py}). Six geometric transforms (identity, horizontal flip,
vertical flip, 90$^\circ$/180$^\circ$/270$^\circ$ rotation) were applied to each
test image, predictions from multiple transforms and multiple nnU-Net checkpoints
were fused via pixel-wise majority voting, and a structured summary was generated.
TTA and ensemble flags could be triggered directly from the natural-language input
(e.g., ``segment teeth with TTA and ensemble'') and were resolved during intent
parsing with automatic dataset defaulting to TDD when unspecified.


% ═══════════════════════════════════════════════════════════════════════════════
% 操作 4 [+]: 新增 Section 2.7
% 插入位置: 2.6 Statistical Analysis 之后，3. Results 之前
% ═══════════════════════════════════════════════════════════════════════════════

\subsection{Intent-to-Workflow Evaluation Protocol}

To systematically assess the agent's decision accuracy, we designed a 6-dimension
weighted evaluation framework (detailed in Section~3.4 and Table~\ref{tab:eval_dims})
and constructed a test set of 11~representative intents spanning four categories:
(i)~standard executable requests (5~intents) for which a validated registry method
existed; (ii)~boundary requests (3~intents) that were in-scope but lacked a
registered method; (iii)~trap requests (2~intents) that were completely out of
dental~CV scope; and (iv)~ambiguous requests (1~intent) with intentionally
underspecified task or dataset information. For each intent, the expected platform
outcome (\texttt{supported}, \texttt{executable}, and target method) was
pre-defined. The platform's actual output was compared against this ground truth,
and a weighted score was computed across the six dimensions.


% =============================================================================
% 最终 .tex 文件结构预览
% =============================================================================
% 2. Materials and methods
%   2.1 Study design and platform overview
%       ... (原文前半保留不变) ...
%       [R] 操作 1: 替换 "For the present study..." 段
%       [+] 操作 2: 新增 "Agent Decision Pipeline" 子节
%   2.2 Datasets and representative workflows
%       ... (原文保留不变) ...
%       [+] 操作 3: 新增 "TTA and Ensemble Configuration" 子节
%   2.3 Workflow execution and automated dataset preparation    ← 原文保留
%   2.4 Validation workflows and outcome measures               ← 原文保留
%   2.5 Skill-based configuration case study                    ← 原文保留
%   2.6 Statistical analysis                                    ← 原文保留
%   2.7 Intent-to-Workflow Evaluation Protocol                  ← [+] 操作 4
% 3. Results
%   3.1 End-to-end workflow execution                           ← 原文保留
%   3.2 Structured reporting                                    ← 原文保留
%   3.3 Expert-informed configuration trials                    ← 原文保留
%   3.4 Intent-to-Workflow Decision Accuracy                    ← [+] section_3_4.tex
% 4. Discussion                                                 ← (待改)
